% GENERATED FILE. Edit canonical lessons, then run npm run generate:books.
% canonical-source-hash: fnv1a64:109d88388fa8d10b
% canonical-lessons: FR-C474-gaiete, FR-C474-rigoler, FR-C474-fou-rire, FR-C474-marrant, FR-C474-chatouiller

\chapter{Rire aux éclats}
\label{ch:fr-rire-aux-eclats}
\hlchaptermodality{474}

\begin{chapteropening}
\textbf{By the end of this chapter:} \emph{I can talk about fun, laughter and good spirits.}

Complete the last lesson of chapter 474: I can talk about fun, laughter and good spirits.
\end{chapteropening}

\section[la gaieté]{la gaieté — cheerfulness}

\label{lesson:FR-C474-gaiete}



\begin{quote}
Before the new one: say the French for modest, then the French for to praise.
\end{quote}

\subsection*{You'll want to know: la gaieté}
\textbf{la gaieté} — \textquotedblleft{}cheerfulness\textquotedblright{}.

\begin{quote}
Sa gaieté met tout le monde de bonne humeur.
\end{quote}

\begin{quote}
\emph{Her cheerfulness puts everyone in a good mood.}
\end{quote}

Also spelt \emph{la gaîté}. Built on \emph{gai} \textquotedblleft{}cheerful\textquotedblright{}.

\begin{cousinweb}
\emph{la gaieté} is built on \emph{gai} \textquotedblleft{}cheerful\textquotedblright{} (EtymDB), which came into French from Old Occitan \emph{gai}; it is usually traced further back to a Germanic word.
\end{cousinweb}

\begin{grammarlens}[title={What you've built}]
One more word for telling, explaining and describing at B1.
\end{grammarlens}

\subsection*{Your turn}
\begin{itemize}
\raggedright
  \item \emph{Say it:} \emph{la gaieté}
  \item \emph{Say it:} \emph{la gaieté}, once more
  \item \emph{Say it:} a sentence of your own with \emph{la gaieté}
  \item \emph{Recall:} say the French for modest, then the French for to praise, then say \emph{la gaieté} again
\end{itemize}

\begin{tcolorbox}[breakable,colback=teal!4,colframe=teal!35!black,title={Before you move on}]
What does la gaieté mean? (\textquotedblleft{}Cheerfulness\textquotedblright{}.) Say it once more.
\end{tcolorbox}
\section[rigoler]{rigoler — to laugh, to have fun (informal)}

\label{lesson:FR-C474-rigoler}



\begin{quote}
Before the new one: say the French for to praise, then the French for cheerfulness.
\end{quote}

\subsection*{You'll want to know: rigoler}
\textbf{rigoler} — \textquotedblleft{}to laugh, to have fun (informal)\textquotedblright{}.

\begin{quote}
On a bien rigolé pendant la soirée.
\end{quote}

\begin{quote}
\emph{We had a good laugh at the party.}
\end{quote}

Informal. \emph{Tu rigoles !} — you are joking!

\begin{cousinweb}
\emph{rigoler} goes back to Old French \emph{rigoler} \textquotedblleft{}to amuse oneself\textquotedblright{}. The origin is not settled; a link with \emph{rire} is often suggested.
\end{cousinweb}

\begin{grammarlens}[title={What you've built}]
One more word for telling, explaining and describing at B1.
\end{grammarlens}

\subsection*{Your turn}
\begin{itemize}
\raggedright
  \item \emph{Say it:} \emph{rigoler}
  \item \emph{Say it:} \emph{rigoler}, once more
  \item \emph{Say it:} a sentence of your own with \emph{rigoler}
  \item \emph{Recall:} say the French for to praise, then the French for cheerfulness, then say \emph{rigoler} again
\end{itemize}

\begin{tcolorbox}[breakable,colback=teal!4,colframe=teal!35!black,title={Before you move on}]
What does rigoler mean? (\textquotedblleft{}To laugh, to have fun (informal)\textquotedblright{}.) Say it once more.
\end{tcolorbox}
\section[le fou rire]{le fou rire — fit of the giggles}

\label{lesson:FR-C474-fou-rire}



\begin{quote}
Before the new one: say the French for cheerfulness, then the French for to laugh.
\end{quote}

\subsection*{You'll want to know: le fou rire}
\textbf{le fou rire} — \textquotedblleft{}fit of the giggles\textquotedblright{}.

\begin{quote}
Pendant le cours, nous avons eu un fou rire.
\end{quote}

\begin{quote}
\emph{During the lesson we got the giggles.}
\end{quote}

\emph{Avoir le fou rire} — to be unable to stop laughing.

\begin{cousinweb}
\emph{le fou rire} is \textquotedblleft{}mad laughter\textquotedblright{}: \emph{fou} is inherited from Latin \emph{follis} \textquotedblleft{}bellows, a bag of wind\textquotedblright{}, and \emph{rire} from Latin \emph{rīdēre} \textquotedblleft{}to laugh\textquotedblright{} (Lewis \& Short).
\end{cousinweb}

\begin{grammarlens}[title={What you've built}]
One more word for telling, explaining and describing at B1.
\end{grammarlens}

\subsection*{Your turn}
\begin{itemize}
\raggedright
  \item \emph{Say it:} \emph{le fou rire}
  \item \emph{Say it:} \emph{le fou rire}, once more
  \item \emph{Say it:} a sentence of your own with \emph{le fou rire}
  \item \emph{Recall:} say the French for cheerfulness, then the French for to laugh, then say \emph{le fou rire} again
\end{itemize}

\begin{tcolorbox}[breakable,colback=teal!4,colframe=teal!35!black,title={Before you move on}]
What does le fou rire mean? (\textquotedblleft{}Fit of the giggles\textquotedblright{}.) Say it once more.
\end{tcolorbox}
\section[marrant]{marrant — funny (informal)}

\label{lesson:FR-C474-marrant}



\begin{quote}
Before the new one: say the French for to laugh, then the French for fit of the giggles.
\end{quote}

\subsection*{You'll want to know: marrant}
\textbf{marrant} — \textquotedblleft{}funny (informal)\textquotedblright{}.

\begin{quote}
Ton frère est vraiment marrant.
\end{quote}

\begin{quote}
\emph{Your brother is really funny.}
\end{quote}

Informal, like \emph{drôle}. \emph{Se marrer} — to laugh.

\begin{cousinweb}
\emph{marrant} is the present participle of \emph{se marrer} \textquotedblleft{}to laugh\textquotedblright{}, a word of familiar French. The origin is not settled.
\end{cousinweb}

\begin{grammarlens}[title={What you've built}]
One more word for telling, explaining and describing at B1.
\end{grammarlens}

\subsection*{Your turn}
\begin{itemize}
\raggedright
  \item \emph{Say it:} \emph{marrant}
  \item \emph{Say it:} \emph{marrant}, once more
  \item \emph{Say it:} a sentence of your own with \emph{marrant}
  \item \emph{Recall:} say the French for to laugh, then the French for fit of the giggles, then say \emph{marrant} again
\end{itemize}

\begin{tcolorbox}[breakable,colback=teal!4,colframe=teal!35!black,title={Before you move on}]
What does marrant mean? (\textquotedblleft{}Funny (informal)\textquotedblright{}.) Say it once more.
\end{tcolorbox}
\section[chatouiller]{chatouiller — to tickle}

\label{lesson:FR-C474-chatouiller}



\begin{quote}
Before the new one: say the French for fit of the giggles, then the French for funny.
\end{quote}

\subsection*{You'll want to know: chatouiller}
\textbf{chatouiller} — \textquotedblleft{}to tickle\textquotedblright{}.

\begin{quote}
Le grand-père chatouille le bébé, qui rit.
\end{quote}

\begin{quote}
\emph{The grandfather tickles the baby, who laughs.}
\end{quote}

\emph{Être chatouilleux} — to be ticklish.

\begin{cousinweb}
The origin of \emph{chatouiller} is not settled. It is probably an expressive word; Germanic words for tickling, like Old English \emph{citelian}, look similar (EtymDB).
\end{cousinweb}

\begin{grammarlens}[title={What you've built}]
One more word for telling, explaining and describing at B1.
\end{grammarlens}

\subsection*{Your turn}
\begin{itemize}
\raggedright
  \item \emph{Say it:} \emph{chatouiller}
  \item \emph{Say it:} \emph{chatouiller}, once more
  \item \emph{Say it:} a sentence of your own with \emph{chatouiller}
  \item \emph{Recall:} say the French for fit of the giggles, then the French for funny, then say \emph{chatouiller} again
\end{itemize}

\begin{tcolorbox}[breakable,colback=teal!4,colframe=teal!35!black,title={Before you move on}]
What does chatouiller mean? (\textquotedblleft{}To tickle\textquotedblright{}.) Say it once more.
\end{tcolorbox}
